\documentclass[10pt, landscape, twocolumn]{article}

% --- Librerías Requeridas ---
\usepackage[margin=0.5in, columnsep=0.4in]{geometry}
\usepackage{amsmath, amssymb, amsthm}
\usepackage{xcolor}
\usepackage{tabularx}
\usepackage{tikz}
\usetikzlibrary{arrows.meta, positioning, calc}
\usepackage{enumitem}
\usepackage{mathpazo} % Tipografía Palatino
\usepackage{booktabs}

% --- Configuración de Formato ---
\pagestyle{empty} % Sin números de página
\setlength{\parindent}{0pt}
\setlength{\parskip}{2pt}

% Ajuste de espacios en ecuaciones
\AtBeginDocument{
  \setlength{\abovedisplayskip}{3pt}
  \setlength{\belowdisplayskip}{3pt}
  \setlength{\abovedisplayshortskip}{1pt}
  \setlength{\belowdisplayshortskip}{1pt}
}

% --- Definición de Colores ---
\definecolor{themeblue}{RGB}{0, 90, 160}
\definecolor{themegray}{RGB}{80, 80, 80}

% --- Comandos Personalizados para Títulos ---
\newcommand{\cheatsheetsection}[1]{
    \vspace{0.6em}
    {\color{themeblue}\large\bfseries #1} \par
    \vspace{-0.5em}
    {\color{themeblue}\rule{\linewidth}{1pt}} \par
    \vspace{0.2em}
}

\newcommand{\cheatsheetsubsection}[1]{
    \vspace{0.4em}
    {\color{themeblue}\normalsize\bfseries #1} \par
    \vspace{0.1em}
}

% --- Inicio del Documento ---
\begin{document}

% --- Cabecera ---
\begin{center}
    \vspace{-1em}
    {\color{themeblue}\huge\bfseries Título de la Plantilla Basal}\\[0.1em]
    {\color{themegray}\small Subtítulo descriptivo o guía rápida para el tema seleccionado}\\[0.4em]
    {\color{themegray}\footnotesize \textbf{Por Tu Nombre Aquí}}
\end{center}
\vspace{-1em}

% --- Sección de Ejemplo ---
\cheatsheetsection{Primera Sección de Ejemplo}
Este es un ejemplo de texto genérico para verificar el comportamiento de la tipografía y los márgenes en las dos columnas. Puedes añadir ecuaciones en línea como $E = mc^2$ o en bloques:
\[ \int_{a}^{b} f(x) \, dx = F(b) - F(a) \]

% --- Subsección de Ejemplo ---
\cheatsheetsubsection{Subsección Detallada}
Aquí puedes incluir listas o descripciones más específicas:
\begin{itemize}[leftmargin=*, nosep]
    \item Primer punto de importancia.
    \item Segundo concepto clave con una nota adicional.
    \item Tercer elemento de la lista.
\end{itemize}

% --- Ejemplo de Tabla ---
\cheatsheetsection{Referencia de Datos (Tabla)}
\vspace{0.2em}
\renewcommand{\arraystretch}{1.2}
\begin{tabularx}{\linewidth}{|p{0.3\linewidth}|X|}
\hline
\textbf{Categoría} & \textbf{Descripción o Aplicación} \\ \hline
Concepto A & Descripción detallada de cómo aplicar este concepto en un entorno práctico. \\ \hline
Concepto B & Nota sobre las limitaciones o excepciones de esta regla específica. \\ \hline
Fórmula C & $\nabla \cdot \mathbf{B} = 0$ - Breve explicación de la ley física. \\ \hline
\end{tabularx}

\vspace{1em}
\cheatsheetsection{Segunda Sección}
Puedes continuar escribiendo aquí. LaTeX balanceará las columnas automáticamente o puedes forzar un salto de columna con el comando \texttt{\textbackslash columnbreak} o empezar una página nueva con \texttt{\textbackslash newpage}.

% --- Espacio para más contenido ---

\end{document}